\chapter{The long tables}
\label{app:tables}

The tables in this appendix are the complete versions of grids the chapters
quote from. Every one is generated by \code{scripts/make\_report\_tables.py} from
a committed result file, and the paths are printed under each.

\section{Every comparison in the family}

\begin{center}
\small
\footnotesize
\begin{longtable}{>{\raggedright\arraybackslash}p{.13\textwidth}>{\raggedright\arraybackslash}p{.11\textwidth}>{\raggedright\arraybackslash}p{.15\textwidth}rrr>{\raggedright\arraybackslash}p{.15\textwidth}}
\toprule
\textbf{Comparison} & \textbf{Metric} & \textbf{Slice} & \textbf{Effect} & \textbf{p} & \textbf{adjusted p} & \textbf{Verdict} \\
\midrule
\endfirsthead
\toprule
\textbf{Comparison} & \textbf{Metric} & \textbf{Slice} & \textbf{Effect} & \textbf{p} & \textbf{adjusted p} & \textbf{Verdict} \\
\midrule
\endhead
\bottomrule
\endfoot
rules over llm & \ident{macro_f1} & \ident{all} & 0.1558 & 0.00781 & 0.02578 & improvement \\  % results: experiments/results/analysis/2026-08-26T21-20-11Z_p6-analysis_6457ac7/analysis.json
rules over llm & \ident{macro_f1} & \ident{seen_locales} & 0.1618 & 0.00586 & 0.02578 & improvement \\  % results: experiments/results/analysis/2026-08-26T21-20-11Z_p6-analysis_6457ac7/analysis.json
rules over llm & \ident{macro_f1} & \ident{unseen_locale} & 0.1231 & 0.05664 & 0.08496 & no significant change \\  % results: experiments/results/analysis/2026-08-26T21-20-11Z_p6-analysis_6457ac7/analysis.json
rules over llm & \ident{macro_f1} & \ident{tier:clean} & 0.1981 & 0.01758 & 0.04143 & improvement \\  % results: experiments/results/analysis/2026-08-26T21-20-11Z_p6-analysis_6457ac7/analysis.json
rules over llm & \ident{macro_f1} & \ident{tier:partial} & 0.2155 & 0.01562 & 0.04143 & improvement \\  % results: experiments/results/analysis/2026-08-26T21-20-11Z_p6-analysis_6457ac7/analysis.json
rules over llm & \ident{macro_f1} & \ident{tier:mixed} & 0.1653 & 0.05664 & 0.08496 & no significant change \\  % results: experiments/results/analysis/2026-08-26T21-20-11Z_p6-analysis_6457ac7/analysis.json
rules over llm & \ident{macro_f1} & \ident{tier:hostile} & -0.0584 & 0.19531 & 0.23872 & no significant change \\  % results: experiments/results/analysis/2026-08-26T21-20-11Z_p6-analysis_6457ac7/analysis.json
rules over llm & \ident{finding_precision} & \ident{MISSING_AUTOCOMPLETE} & 0.2197 & 0.00781 & 0.02578 & improvement \\  % results: experiments/results/analysis/2026-08-26T21-20-11Z_p6-analysis_6457ac7/analysis.json
rules over llm & \ident{finding_recall} & \ident{MISSING_AUTOCOMPLETE} & -0.1152 & 0.00195 & 0.02578 & regression \\  % results: experiments/results/analysis/2026-08-26T21-20-11Z_p6-analysis_6457ac7/analysis.json
rules over llm & \ident{finding_precision} & \ident{WRONG_AUTOCOMPLETE} & 0.7680 & 0.00781 & 0.02578 & improvement \\  % results: experiments/results/analysis/2026-08-26T21-20-11Z_p6-analysis_6457ac7/analysis.json
rules over llm & \ident{finding_recall} & \ident{WRONG_AUTOCOMPLETE} & -0.0333 & 1.00000 & 1.00000 & no significant change \\  % results: experiments/results/analysis/2026-08-26T21-20-11Z_p6-analysis_6457ac7/analysis.json
ngram over llm & \ident{macro_f1} & \ident{all} & 0.1041 & 0.08594 & 0.11816 & no significant change \\  % results: experiments/results/analysis/2026-08-26T21-20-11Z_p6-analysis_6457ac7/analysis.json
ngram over llm & \ident{macro_f1} & \ident{seen_locales} & 0.1362 & 0.03516 & 0.06445 & no significant change \\  % results: experiments/results/analysis/2026-08-26T21-20-11Z_p6-analysis_6457ac7/analysis.json
ngram over llm & \ident{macro_f1} & \ident{unseen_locale} & -0.0409 & 0.37500 & 0.41250 & no significant change \\  % results: experiments/results/analysis/2026-08-26T21-20-11Z_p6-analysis_6457ac7/analysis.json
ngram over llm & \ident{macro_f1} & \ident{tier:clean} & 0.2654 & 0.00391 & 0.02578 & improvement \\  % results: experiments/results/analysis/2026-08-26T21-20-11Z_p6-analysis_6457ac7/analysis.json
ngram over llm & \ident{macro_f1} & \ident{tier:partial} & 0.2768 & 0.00391 & 0.02578 & improvement \\  % results: experiments/results/analysis/2026-08-26T21-20-11Z_p6-analysis_6457ac7/analysis.json
ngram over llm & \ident{macro_f1} & \ident{tier:mixed} & 0.0799 & 0.30664 & 0.34894 & no significant change \\  % results: experiments/results/analysis/2026-08-26T21-20-11Z_p6-analysis_6457ac7/analysis.json
ngram over llm & \ident{macro_f1} & \ident{tier:hostile} & 0.0143 & 0.79102 & 0.81573 & no significant change \\  % results: experiments/results/analysis/2026-08-26T21-20-11Z_p6-analysis_6457ac7/analysis.json
ngram over llm & \ident{finding_precision} & \ident{MISSING_AUTOCOMPLETE} & 0.1927 & 0.00781 & 0.02578 & improvement \\  % results: experiments/results/analysis/2026-08-26T21-20-11Z_p6-analysis_6457ac7/analysis.json
ngram over llm & \ident{finding_recall} & \ident{MISSING_AUTOCOMPLETE} & -0.2722 & 0.00195 & 0.02578 & regression \\  % results: experiments/results/analysis/2026-08-26T21-20-11Z_p6-analysis_6457ac7/analysis.json
ngram over llm & \ident{finding_precision} & \ident{WRONG_AUTOCOMPLETE} & 0.6810 & 0.00781 & 0.02578 & improvement \\  % results: experiments/results/analysis/2026-08-26T21-20-11Z_p6-analysis_6457ac7/analysis.json
ngram over llm & \ident{finding_recall} & \ident{WRONG_AUTOCOMPLETE} & -0.2667 & 0.03125 & 0.06066 & no significant change \\  % results: experiments/results/analysis/2026-08-26T21-20-11Z_p6-analysis_6457ac7/analysis.json
ngram over rules & \ident{macro_f1} & \ident{all} & -0.0517 & 0.22656 & 0.26702 & no significant change \\  % results: experiments/results/analysis/2026-08-26T21-20-11Z_p6-analysis_6457ac7/analysis.json
ngram over rules & \ident{macro_f1} & \ident{seen_locales} & -0.0256 & 0.59375 & 0.63206 & no significant change \\  % results: experiments/results/analysis/2026-08-26T21-20-11Z_p6-analysis_6457ac7/analysis.json
ngram over rules & \ident{macro_f1} & \ident{unseen_locale} & -0.1640 & 0.02148 & 0.04727 & regression \\  % results: experiments/results/analysis/2026-08-26T21-20-11Z_p6-analysis_6457ac7/analysis.json
ngram over rules & \ident{macro_f1} & \ident{tier:clean} & 0.0673 & 0.02344 & 0.04834 & improvement \\  % results: experiments/results/analysis/2026-08-26T21-20-11Z_p6-analysis_6457ac7/analysis.json
ngram over rules & \ident{macro_f1} & \ident{tier:partial} & 0.0613 & 0.05469 & 0.08496 & no significant change \\  % results: experiments/results/analysis/2026-08-26T21-20-11Z_p6-analysis_6457ac7/analysis.json
ngram over rules & \ident{macro_f1} & \ident{tier:mixed} & -0.0854 & 0.15820 & 0.20080 & no significant change \\  % results: experiments/results/analysis/2026-08-26T21-20-11Z_p6-analysis_6457ac7/analysis.json
ngram over rules & \ident{macro_f1} & \ident{tier:hostile} & 0.0727 & 0.08203 & 0.11770 & no significant change \\  % results: experiments/results/analysis/2026-08-26T21-20-11Z_p6-analysis_6457ac7/analysis.json
ngram over rules & \ident{finding_precision} & \ident{MISSING_AUTOCOMPLETE} & -0.0269 & 0.01758 & 0.04143 & regression \\  % results: experiments/results/analysis/2026-08-26T21-20-11Z_p6-analysis_6457ac7/analysis.json
ngram over rules & \ident{finding_recall} & \ident{MISSING_AUTOCOMPLETE} & -0.1570 & 0.01758 & 0.04143 & regression \\  % results: experiments/results/analysis/2026-08-26T21-20-11Z_p6-analysis_6457ac7/analysis.json
ngram over rules & \ident{finding_precision} & \ident{WRONG_AUTOCOMPLETE} & -0.0870 & 0.04688 & 0.08141 & no significant change \\  % results: experiments/results/analysis/2026-08-26T21-20-11Z_p6-analysis_6457ac7/analysis.json
ngram over rules & \ident{finding_recall} & \ident{WRONG_AUTOCOMPLETE} & -0.2333 & 0.09375 & 0.12375 & no significant change \\  % results: experiments/results/analysis/2026-08-26T21-20-11Z_p6-analysis_6457ac7/analysis.json
\end{longtable}
\par\smallskip\noindent{\footnotesize Source: \rpath{experiments/results/analysis/2026-08-26T21-20-11Z_p6-analysis_6457ac7/analysis.json}}

\captionof{table}{All comparisons in the corrected family, read in the direction
they were measured. The \emph{Comparison} column is built from the baseline, the
candidate and the effect rather than from the comparison name, which reads
backwards at a glance.}
\label{tab:comparisons}
\end{center}

\section{The locale by tier grid}

\begin{center}
\small
\begin{longtable}{llrrrr}
\toprule
\textbf{Locale} & \textbf{Tier} & \textbf{Fields} & \textbf{rules} & \textbf{ngram} & \textbf{llm} \\
\midrule
\endfirsthead
\toprule
\textbf{Locale} & \textbf{Tier} & \textbf{Fields} & \textbf{rules} & \textbf{ngram} & \textbf{llm} \\
\midrule
\endhead
\bottomrule
\endfoot
de-DE & clean & 84 & 0.7962 & 0.9612 & 0.6732 \\  % results: experiments/results/bench/2026-08-26T20-32-48Z_bench_6457ac7/benchmark.json
de-DE & hostile & 104 & 0.3838 & 0.5401 & 0.4420 \\  % results: experiments/results/bench/2026-08-26T20-32-48Z_bench_6457ac7/benchmark.json
de-DE & mixed & 104 & 0.7518 & 0.7561 & 0.5888 \\  % results: experiments/results/bench/2026-08-26T20-32-48Z_bench_6457ac7/benchmark.json
de-DE & partial & 84 & 0.7962 & 0.9522 & 0.6740 \\  % results: experiments/results/bench/2026-08-26T20-32-48Z_bench_6457ac7/benchmark.json
en-GB & clean & 88 & 0.8741 & 0.9902 & 0.6029 \\  % results: experiments/results/bench/2026-08-26T20-32-48Z_bench_6457ac7/benchmark.json
en-GB & hostile & 108 & 0.4667 & 0.6258 & 0.5611 \\  % results: experiments/results/bench/2026-08-26T20-32-48Z_bench_6457ac7/benchmark.json
en-GB & mixed & 108 & 0.7493 & 0.7829 & 0.6478 \\  % results: experiments/results/bench/2026-08-26T20-32-48Z_bench_6457ac7/benchmark.json
en-GB & partial & 88 & 0.8741 & 0.9902 & 0.6826 \\  % results: experiments/results/bench/2026-08-26T20-32-48Z_bench_6457ac7/benchmark.json
en-NG & clean & 85 & 0.8741 & 0.9902 & 0.7011 \\  % results: experiments/results/bench/2026-08-26T20-32-48Z_bench_6457ac7/benchmark.json
en-NG & hostile & 104 & 0.4453 & 0.6284 & 0.5656 \\  % results: experiments/results/bench/2026-08-26T20-32-48Z_bench_6457ac7/benchmark.json
en-NG & mixed & 104 & 0.7790 & 0.8753 & 0.6555 \\  % results: experiments/results/bench/2026-08-26T20-32-48Z_bench_6457ac7/benchmark.json
en-NG & partial & 84 & 0.8741 & 0.9902 & 0.6791 \\  % results: experiments/results/bench/2026-08-26T20-32-48Z_bench_6457ac7/benchmark.json
en-US & clean & 88 & 0.8741 & 0.9902 & 0.7210 \\  % results: experiments/results/bench/2026-08-26T20-32-48Z_bench_6457ac7/benchmark.json
en-US & hostile & 108 & 0.4667 & 0.5704 & 0.5333 \\  % results: experiments/results/bench/2026-08-26T20-32-48Z_bench_6457ac7/benchmark.json
en-US & mixed & 108 & 0.8068 & 0.7727 & 0.6267 \\  % results: experiments/results/bench/2026-08-26T20-32-48Z_bench_6457ac7/benchmark.json
en-US & partial & 88 & 0.8741 & 0.9902 & 0.6251 \\  % results: experiments/results/bench/2026-08-26T20-32-48Z_bench_6457ac7/benchmark.json
fr-FR & clean & 84 & 0.7962 & 0.6919 & 0.6373 \\  % results: experiments/results/bench/2026-08-26T20-32-48Z_bench_6457ac7/benchmark.json
fr-FR & hostile & 104 & 0.4002 & 0.3347 & 0.4733 \\  % results: experiments/results/bench/2026-08-26T20-32-48Z_bench_6457ac7/benchmark.json
fr-FR & mixed & 104 & 0.7366 & 0.5842 & 0.6330 \\  % results: experiments/results/bench/2026-08-26T20-32-48Z_bench_6457ac7/benchmark.json
fr-FR & partial & 84 & 0.7962 & 0.6919 & 0.5720 \\  % results: experiments/results/bench/2026-08-26T20-32-48Z_bench_6457ac7/benchmark.json
ja-JP & clean & 91 & 0.9074 & 0.9349 & 0.6100 \\  % results: experiments/results/bench/2026-08-26T20-32-48Z_bench_6457ac7/benchmark.json
ja-JP & hostile & 111 & 0.5060 & 0.5125 & 0.4917 \\  % results: experiments/results/bench/2026-08-26T20-32-48Z_bench_6457ac7/benchmark.json
ja-JP & mixed & 111 & 0.7937 & 0.7098 & 0.6259 \\  % results: experiments/results/bench/2026-08-26T20-32-48Z_bench_6457ac7/benchmark.json
ja-JP & partial & 91 & 0.9074 & 0.9016 & 0.6338 \\  % results: experiments/results/bench/2026-08-26T20-32-48Z_bench_6457ac7/benchmark.json
\end{longtable}
\par\smallskip\noindent{\footnotesize Source: \rpath{experiments/results/bench/2026-08-26T20-32-48Z_bench_6457ac7/benchmark.json}}

\captionof{table}{Macro-F1 for every locale and tier cell, all three engines.
Every cell clears the reporting minimum, so nothing here is suppressed.}
\label{tab:locale-by-tier}
\end{center}

\section{Per label behaviour}

\begin{center}
\small
\begin{longtable}{lrrrr}
\toprule
\textbf{Label} & \textbf{Support} & \textbf{rules} & \textbf{ngram} & \textbf{llm} \\
\midrule
\endfirsthead
\toprule
\textbf{Label} & \textbf{Support} & \textbf{rules} & \textbf{ngram} & \textbf{llm} \\
\midrule
\endhead
\bottomrule
\endfoot
\texttt{CC\_EXP\_SPLIT\_MONTH} & 48 & 1.0000 & 1.0000 & 0.6047 \\  % results: experiments/results/test/2026-08-26T21-15-35Z_p6-rules_6457ac7/metrics.json experiments/results/test/2026-08-26T21-17-53Z_p6-ngram_6457ac7/metrics.json experiments/results/test/2026-08-26T20-32-48Z_p6-llm_6457ac7/metrics.json
\texttt{CC\_EXP\_SPLIT\_YEAR} & 48 & 1.0000 & 1.0000 & 0.6842 \\  % results: experiments/results/test/2026-08-26T21-15-35Z_p6-rules_6457ac7/metrics.json experiments/results/test/2026-08-26T21-17-53Z_p6-ngram_6457ac7/metrics.json experiments/results/test/2026-08-26T20-32-48Z_p6-llm_6457ac7/metrics.json
\texttt{COMPOSITE\_UNSPLIT} & 48 & 0.8293 & 0.9565 & 0.0000 \\  % results: experiments/results/test/2026-08-26T21-15-35Z_p6-rules_6457ac7/metrics.json experiments/results/test/2026-08-26T21-17-53Z_p6-ngram_6457ac7/metrics.json experiments/results/test/2026-08-26T20-32-48Z_p6-llm_6457ac7/metrics.json
\texttt{NOT\_AUTOFILLABLE} & 336 & 0.8819 & 0.9587 & 0.7147 \\  % results: experiments/results/test/2026-08-26T21-15-35Z_p6-rules_6457ac7/metrics.json experiments/results/test/2026-08-26T21-17-53Z_p6-ngram_6457ac7/metrics.json experiments/results/test/2026-08-26T20-32-48Z_p6-llm_6457ac7/metrics.json
\texttt{UNKNOWN} & 120 & 0.3301 & 0.7255 & 0.2458 \\  % results: experiments/results/test/2026-08-26T21-15-35Z_p6-rules_6457ac7/metrics.json experiments/results/test/2026-08-26T21-17-53Z_p6-ngram_6457ac7/metrics.json experiments/results/test/2026-08-26T20-32-48Z_p6-llm_6457ac7/metrics.json
\texttt{address-level1} & 64 & 1.0000 & 0.7211 & 0.9375 \\  % results: experiments/results/test/2026-08-26T21-15-35Z_p6-rules_6457ac7/metrics.json experiments/results/test/2026-08-26T21-17-53Z_p6-ngram_6457ac7/metrics.json experiments/results/test/2026-08-26T20-32-48Z_p6-llm_6457ac7/metrics.json
\texttt{address-level2} & 84 & 0.8000 & 0.6900 & 0.8299 \\  % results: experiments/results/test/2026-08-26T21-15-35Z_p6-rules_6457ac7/metrics.json experiments/results/test/2026-08-26T21-17-53Z_p6-ngram_6457ac7/metrics.json experiments/results/test/2026-08-26T20-32-48Z_p6-llm_6457ac7/metrics.json
\texttt{address-line1} & 24 & 0.6286 & 0.6032 & 0.2901 \\  % results: experiments/results/test/2026-08-26T21-15-35Z_p6-rules_6457ac7/metrics.json experiments/results/test/2026-08-26T21-17-53Z_p6-ngram_6457ac7/metrics.json experiments/results/test/2026-08-26T20-32-48Z_p6-llm_6457ac7/metrics.json
\texttt{address-line2} & 24 & 0.9091 & 0.6364 & 0.7407 \\  % results: experiments/results/test/2026-08-26T21-15-35Z_p6-rules_6457ac7/metrics.json experiments/results/test/2026-08-26T21-17-53Z_p6-ngram_6457ac7/metrics.json experiments/results/test/2026-08-26T20-32-48Z_p6-llm_6457ac7/metrics.json
\texttt{address-line3} & 24 & 0.8837 & 0.6842 & 0.7917 \\  % results: experiments/results/test/2026-08-26T21-15-35Z_p6-rules_6457ac7/metrics.json experiments/results/test/2026-08-26T21-17-53Z_p6-ngram_6457ac7/metrics.json experiments/results/test/2026-08-26T20-32-48Z_p6-llm_6457ac7/metrics.json
\texttt{bday} & 24 & 0.8293 & 0.9362 & 0.8293 \\  % results: experiments/results/test/2026-08-26T21-15-35Z_p6-rules_6457ac7/metrics.json experiments/results/test/2026-08-26T21-17-53Z_p6-ngram_6457ac7/metrics.json experiments/results/test/2026-08-26T20-32-48Z_p6-llm_6457ac7/metrics.json
\texttt{cc-csc} & 96 & 0.8000 & 0.8804 & 0.8977 \\  % results: experiments/results/test/2026-08-26T21-15-35Z_p6-rules_6457ac7/metrics.json experiments/results/test/2026-08-26T21-17-53Z_p6-ngram_6457ac7/metrics.json experiments/results/test/2026-08-26T20-32-48Z_p6-llm_6457ac7/metrics.json
\texttt{cc-exp} & 24 & 0.9333 & 0.8837 & 0.5882 \\  % results: experiments/results/test/2026-08-26T21-15-35Z_p6-rules_6457ac7/metrics.json experiments/results/test/2026-08-26T21-17-53Z_p6-ngram_6457ac7/metrics.json experiments/results/test/2026-08-26T20-32-48Z_p6-llm_6457ac7/metrics.json
\texttt{cc-exp-month} & 24 & 1.0000 & 0.8235 & 0.6111 \\  % results: experiments/results/test/2026-08-26T21-15-35Z_p6-rules_6457ac7/metrics.json experiments/results/test/2026-08-26T21-17-53Z_p6-ngram_6457ac7/metrics.json experiments/results/test/2026-08-26T20-32-48Z_p6-llm_6457ac7/metrics.json
\texttt{cc-exp-year} & 24 & 0.7692 & 0.6667 & 0.6111 \\  % results: experiments/results/test/2026-08-26T21-15-35Z_p6-rules_6457ac7/metrics.json experiments/results/test/2026-08-26T21-17-53Z_p6-ngram_6457ac7/metrics.json experiments/results/test/2026-08-26T20-32-48Z_p6-llm_6457ac7/metrics.json
\texttt{cc-name} & 48 & 0.9677 & 0.7547 & 0.8276 \\  % results: experiments/results/test/2026-08-26T21-15-35Z_p6-rules_6457ac7/metrics.json experiments/results/test/2026-08-26T21-17-53Z_p6-ngram_6457ac7/metrics.json experiments/results/test/2026-08-26T20-32-48Z_p6-llm_6457ac7/metrics.json
\texttt{cc-number} & 96 & 0.8075 & 0.8639 & 0.9005 \\  % results: experiments/results/test/2026-08-26T21-15-35Z_p6-rules_6457ac7/metrics.json experiments/results/test/2026-08-26T21-17-53Z_p6-ngram_6457ac7/metrics.json experiments/results/test/2026-08-26T20-32-48Z_p6-llm_6457ac7/metrics.json
\texttt{cc-type} & 72 & 1.0000 & 0.9714 & 0.5743 \\  % results: experiments/results/test/2026-08-26T21-15-35Z_p6-rules_6457ac7/metrics.json experiments/results/test/2026-08-26T21-17-53Z_p6-ngram_6457ac7/metrics.json experiments/results/test/2026-08-26T20-32-48Z_p6-llm_6457ac7/metrics.json
\texttt{country} & 96 & 0.8421 & 0.7630 & 0.3051 \\  % results: experiments/results/test/2026-08-26T21-15-35Z_p6-rules_6457ac7/metrics.json experiments/results/test/2026-08-26T21-17-53Z_p6-ngram_6457ac7/metrics.json experiments/results/test/2026-08-26T20-32-48Z_p6-llm_6457ac7/metrics.json
\texttt{country-name} & 48 & 0.2545 & 0.6250 & 0.4458 \\  % results: experiments/results/test/2026-08-26T21-15-35Z_p6-rules_6457ac7/metrics.json experiments/results/test/2026-08-26T21-17-53Z_p6-ngram_6457ac7/metrics.json experiments/results/test/2026-08-26T20-32-48Z_p6-llm_6457ac7/metrics.json
\texttt{current-password} & 48 & 0.0000 & 0.8889 & 0.6667 \\  % results: experiments/results/test/2026-08-26T21-15-35Z_p6-rules_6457ac7/metrics.json experiments/results/test/2026-08-26T21-17-53Z_p6-ngram_6457ac7/metrics.json experiments/results/test/2026-08-26T20-32-48Z_p6-llm_6457ac7/metrics.json
\texttt{email} & 144 & 0.9565 & 0.9329 & 0.9531 \\  % results: experiments/results/test/2026-08-26T21-15-35Z_p6-rules_6457ac7/metrics.json experiments/results/test/2026-08-26T21-17-53Z_p6-ngram_6457ac7/metrics.json experiments/results/test/2026-08-26T20-32-48Z_p6-llm_6457ac7/metrics.json
\texttt{family-name} & 84 & 0.9091 & 0.7547 & 0.7891 \\  % results: experiments/results/test/2026-08-26T21-15-35Z_p6-rules_6457ac7/metrics.json experiments/results/test/2026-08-26T21-17-53Z_p6-ngram_6457ac7/metrics.json experiments/results/test/2026-08-26T20-32-48Z_p6-llm_6457ac7/metrics.json
\texttt{given-name} & 84 & 0.9091 & 0.7027 & 0.7750 \\  % results: experiments/results/test/2026-08-26T21-15-35Z_p6-rules_6457ac7/metrics.json experiments/results/test/2026-08-26T21-17-53Z_p6-ngram_6457ac7/metrics.json experiments/results/test/2026-08-26T20-32-48Z_p6-llm_6457ac7/metrics.json
\texttt{name} & 24 & 1.0000 & 0.8163 & 0.7164 \\  % results: experiments/results/test/2026-08-26T21-15-35Z_p6-rules_6457ac7/metrics.json experiments/results/test/2026-08-26T21-17-53Z_p6-ngram_6457ac7/metrics.json experiments/results/test/2026-08-26T20-32-48Z_p6-llm_6457ac7/metrics.json
\texttt{new-password} & 96 & 0.7114 & 0.9495 & 0.8108 \\  % results: experiments/results/test/2026-08-26T21-15-35Z_p6-rules_6457ac7/metrics.json experiments/results/test/2026-08-26T21-17-53Z_p6-ngram_6457ac7/metrics.json experiments/results/test/2026-08-26T20-32-48Z_p6-llm_6457ac7/metrics.json
\texttt{one-time-code} & 48 & 0.7692 & 0.6917 & 0.7105 \\  % results: experiments/results/test/2026-08-26T21-15-35Z_p6-rules_6457ac7/metrics.json experiments/results/test/2026-08-26T21-17-53Z_p6-ngram_6457ac7/metrics.json experiments/results/test/2026-08-26T20-32-48Z_p6-llm_6457ac7/metrics.json
\texttt{organization} & 48 & 1.0000 & 0.7273 & 0.7532 \\  % results: experiments/results/test/2026-08-26T21-15-35Z_p6-rules_6457ac7/metrics.json experiments/results/test/2026-08-26T21-17-53Z_p6-ngram_6457ac7/metrics.json experiments/results/test/2026-08-26T20-32-48Z_p6-llm_6457ac7/metrics.json
\texttt{postal-code} & 105 & 0.7861 & 0.8542 & 0.8646 \\  % results: experiments/results/test/2026-08-26T21-15-35Z_p6-rules_6457ac7/metrics.json experiments/results/test/2026-08-26T21-17-53Z_p6-ngram_6457ac7/metrics.json experiments/results/test/2026-08-26T20-32-48Z_p6-llm_6457ac7/metrics.json
\texttt{sex} & 24 & 1.0000 & 0.9020 & 0.5882 \\  % results: experiments/results/test/2026-08-26T21-15-35Z_p6-rules_6457ac7/metrics.json experiments/results/test/2026-08-26T21-17-53Z_p6-ngram_6457ac7/metrics.json experiments/results/test/2026-08-26T20-32-48Z_p6-llm_6457ac7/metrics.json
\texttt{street-address} & 72 & 0.6154 & 0.6142 & 0.3409 \\  % results: experiments/results/test/2026-08-26T21-15-35Z_p6-rules_6457ac7/metrics.json experiments/results/test/2026-08-26T21-17-53Z_p6-ngram_6457ac7/metrics.json experiments/results/test/2026-08-26T20-32-48Z_p6-llm_6457ac7/metrics.json
\texttt{tel} & 24 & 0.6071 & 0.7200 & 0.6102 \\  % results: experiments/results/test/2026-08-26T21-15-35Z_p6-rules_6457ac7/metrics.json experiments/results/test/2026-08-26T21-17-53Z_p6-ngram_6457ac7/metrics.json experiments/results/test/2026-08-26T20-32-48Z_p6-llm_6457ac7/metrics.json
\texttt{tel-country-code} & 24 & 1.0000 & 1.0000 & 0.8000 \\  % results: experiments/results/test/2026-08-26T21-15-35Z_p6-rules_6457ac7/metrics.json experiments/results/test/2026-08-26T21-17-53Z_p6-ngram_6457ac7/metrics.json experiments/results/test/2026-08-26T20-32-48Z_p6-llm_6457ac7/metrics.json
\texttt{tel-national} & 24 & 0.0000 & 0.5862 & 0.5882 \\  % results: experiments/results/test/2026-08-26T21-15-35Z_p6-rules_6457ac7/metrics.json experiments/results/test/2026-08-26T21-17-53Z_p6-ngram_6457ac7/metrics.json experiments/results/test/2026-08-26T20-32-48Z_p6-llm_6457ac7/metrics.json
\texttt{url} & 24 & 0.8571 & 0.8571 & 0.8571 \\  % results: experiments/results/test/2026-08-26T21-15-35Z_p6-rules_6457ac7/metrics.json experiments/results/test/2026-08-26T21-17-53Z_p6-ngram_6457ac7/metrics.json experiments/results/test/2026-08-26T20-32-48Z_p6-llm_6457ac7/metrics.json
\texttt{username} & 72 & 0.8923 & 0.6949 & 0.8872 \\  % results: experiments/results/test/2026-08-26T21-15-35Z_p6-rules_6457ac7/metrics.json experiments/results/test/2026-08-26T21-17-53Z_p6-ngram_6457ac7/metrics.json experiments/results/test/2026-08-26T20-32-48Z_p6-llm_6457ac7/metrics.json
\end{longtable}
\par\noindent{\footnotesize Source: \rpath{experiments/results/test/2026-08-26T21-15-35Z_p6-rules_6457ac7/metrics.json}}
\par\noindent{\footnotesize Source: \rpath{experiments/results/test/2026-08-26T21-17-53Z_p6-ngram_6457ac7/metrics.json}}
\par\noindent{\footnotesize Source: \rpath{experiments/results/test/2026-08-26T20-32-48Z_p6-llm_6457ac7/metrics.json}}

\captionof{table}{F1 per observed label per engine, with the test split support
beside it. Read the rows against the support column: several labels are carried
by one template in one position.}
\label{tab:per-label}
\end{center}

\section{Confusions}

\begin{center}
\small
\begingroup\small\setlength{\tabcolsep}{4pt}
\begin{tabular}{lllr}
\toprule
\textbf{Engine} & \textbf{Answer key} & \textbf{Predicted} & \textbf{Fields} \\
\midrule
\texttt{rules} & \texttt{NOT\_AUTOFILLABLE} & \texttt{UNKNOWN} & 71 \\  % results: experiments/results/test/2026-08-26T21-15-35Z_p6-rules_6457ac7/metrics.json
 & \texttt{current-password} & \texttt{UNKNOWN} & 48 \\  % results: experiments/results/test/2026-08-26T21-15-35Z_p6-rules_6457ac7/metrics.json
 & \texttt{new-password} & \texttt{UNKNOWN} & 43 \\  % results: experiments/results/test/2026-08-26T21-15-35Z_p6-rules_6457ac7/metrics.json
 & \texttt{street-address} & \texttt{UNKNOWN} & 40 \\  % results: experiments/results/test/2026-08-26T21-15-35Z_p6-rules_6457ac7/metrics.json
 & \texttt{postal-code} & \texttt{UNKNOWN} & 37 \\  % results: experiments/results/test/2026-08-26T21-15-35Z_p6-rules_6457ac7/metrics.json
 & \texttt{country-name} & \texttt{country} & 36 \\  % results: experiments/results/test/2026-08-26T21-15-35Z_p6-rules_6457ac7/metrics.json
 & \texttt{cc-csc} & \texttt{UNKNOWN} & 32 \\  % results: experiments/results/test/2026-08-26T21-15-35Z_p6-rules_6457ac7/metrics.json
 & \texttt{cc-number} & \texttt{UNKNOWN} & 31 \\  % results: experiments/results/test/2026-08-26T21-15-35Z_p6-rules_6457ac7/metrics.json
 & \texttt{address-level2} & \texttt{UNKNOWN} & 28 \\  % results: experiments/results/test/2026-08-26T21-15-35Z_p6-rules_6457ac7/metrics.json
 & \texttt{one-time-code} & \texttt{UNKNOWN} & 18 \\  % results: experiments/results/test/2026-08-26T21-15-35Z_p6-rules_6457ac7/metrics.json
\midrule
\texttt{ngram} & \texttt{country} & \texttt{address-level1} & 30 \\  % results: experiments/results/test/2026-08-26T21-17-53Z_p6-ngram_6457ac7/metrics.json
 & \texttt{username} & \texttt{given-name} & 17 \\  % results: experiments/results/test/2026-08-26T21-17-53Z_p6-ngram_6457ac7/metrics.json
 & \texttt{cc-number} & \texttt{cc-name} & 13 \\  % results: experiments/results/test/2026-08-26T21-17-53Z_p6-ngram_6457ac7/metrics.json
 & \texttt{street-address} & \texttt{one-time-code} & 12 \\  % results: experiments/results/test/2026-08-26T21-17-53Z_p6-ngram_6457ac7/metrics.json
 & \texttt{address-level1} & \texttt{country} & 11 \\  % results: experiments/results/test/2026-08-26T21-17-53Z_p6-ngram_6457ac7/metrics.json
 & \texttt{UNKNOWN} & \texttt{one-time-code} & 10 \\  % results: experiments/results/test/2026-08-26T21-17-53Z_p6-ngram_6457ac7/metrics.json
 & \texttt{UNKNOWN} & \texttt{address-line2} & 9 \\  % results: experiments/results/test/2026-08-26T21-17-53Z_p6-ngram_6457ac7/metrics.json
 & \texttt{country-name} & \texttt{address-level2} & 9 \\  % results: experiments/results/test/2026-08-26T21-17-53Z_p6-ngram_6457ac7/metrics.json
 & \texttt{current-password} & \texttt{new-password} & 8 \\  % results: experiments/results/test/2026-08-26T21-17-53Z_p6-ngram_6457ac7/metrics.json
 & \texttt{NOT\_AUTOFILLABLE} & \texttt{email} & 7 \\  % results: experiments/results/test/2026-08-26T21-17-53Z_p6-ngram_6457ac7/metrics.json
\midrule
\texttt{llm} & \texttt{UNKNOWN} & \texttt{NOT\_AUTOFILLABLE} & 91 \\  % results: experiments/results/test/2026-08-26T20-32-48Z_p6-llm_6457ac7/metrics.json
 & \texttt{country} & \texttt{country-name} & 76 \\  % results: experiments/results/test/2026-08-26T20-32-48Z_p6-llm_6457ac7/metrics.json
 & \texttt{cc-type} & \texttt{NOT\_AUTOFILLABLE} & 43 \\  % results: experiments/results/test/2026-08-26T20-32-48Z_p6-llm_6457ac7/metrics.json
 & \texttt{street-address} & \texttt{address-line1} & 40 \\  % results: experiments/results/test/2026-08-26T20-32-48Z_p6-llm_6457ac7/metrics.json
 & \texttt{COMPOSITE\_UNSPLIT} & \texttt{address-line1} & 36 \\  % results: experiments/results/test/2026-08-26T20-32-48Z_p6-llm_6457ac7/metrics.json
 & \texttt{CC\_EXP\_SPLIT\_MONTH} & \texttt{cc-exp-month} & 22 \\  % results: experiments/results/test/2026-08-26T20-32-48Z_p6-llm_6457ac7/metrics.json
 & \texttt{CC\_EXP\_SPLIT\_YEAR} & \texttt{cc-exp-year} & 22 \\  % results: experiments/results/test/2026-08-26T20-32-48Z_p6-llm_6457ac7/metrics.json
 & \texttt{new-password} & \texttt{current-password} & 20 \\  % results: experiments/results/test/2026-08-26T20-32-48Z_p6-llm_6457ac7/metrics.json
 & \texttt{one-time-code} & \texttt{NOT\_AUTOFILLABLE} & 15 \\  % results: experiments/results/test/2026-08-26T20-32-48Z_p6-llm_6457ac7/metrics.json
 & \texttt{current-password} & \texttt{new-password} & 14 \\  % results: experiments/results/test/2026-08-26T20-32-48Z_p6-llm_6457ac7/metrics.json
\bottomrule
\end{tabular}
\endgroup
\par\noindent{\footnotesize Source: \rpath{experiments/results/test/2026-08-26T21-15-35Z_p6-rules_6457ac7/metrics.json}}
\par\noindent{\footnotesize Source: \rpath{experiments/results/test/2026-08-26T21-17-53Z_p6-ngram_6457ac7/metrics.json}}
\par\noindent{\footnotesize Source: \rpath{experiments/results/test/2026-08-26T20-32-48Z_p6-llm_6457ac7/metrics.json}}

\captionof{table}{The largest confusions each engine makes on the test split.}
\label{tab:confusions}
\end{center}

\begin{center}
\small
\begingroup\small\setlength{\tabcolsep}{4pt}
\begin{tabular}{lrrr}
\toprule
\textbf{Named pair} & \textbf{rules} & \textbf{ngram} & \textbf{llm} \\
\midrule
\texttt{address-level1 as address-level2} & 0 & 0 & 0 \\  % results: experiments/results/test/2026-08-26T21-15-35Z_p6-rules_6457ac7/metrics.json experiments/results/test/2026-08-26T21-17-53Z_p6-ngram_6457ac7/metrics.json experiments/results/test/2026-08-26T20-32-48Z_p6-llm_6457ac7/metrics.json
\texttt{address-level2 as address-level1} & 0 & 0 & 2 \\  % results: experiments/results/test/2026-08-26T21-15-35Z_p6-rules_6457ac7/metrics.json experiments/results/test/2026-08-26T21-17-53Z_p6-ngram_6457ac7/metrics.json experiments/results/test/2026-08-26T20-32-48Z_p6-llm_6457ac7/metrics.json
\texttt{cc-exp as cc-exp-month} & 0 & 1 & 4 \\  % results: experiments/results/test/2026-08-26T21-15-35Z_p6-rules_6457ac7/metrics.json experiments/results/test/2026-08-26T21-17-53Z_p6-ngram_6457ac7/metrics.json experiments/results/test/2026-08-26T20-32-48Z_p6-llm_6457ac7/metrics.json
\texttt{cc-exp-month as cc-exp} & 0 & 0 & 0 \\  % results: experiments/results/test/2026-08-26T21-15-35Z_p6-rules_6457ac7/metrics.json experiments/results/test/2026-08-26T21-17-53Z_p6-ngram_6457ac7/metrics.json experiments/results/test/2026-08-26T20-32-48Z_p6-llm_6457ac7/metrics.json
\texttt{email as username} & 0 & 0 & 0 \\  % results: experiments/results/test/2026-08-26T21-15-35Z_p6-rules_6457ac7/metrics.json experiments/results/test/2026-08-26T21-17-53Z_p6-ngram_6457ac7/metrics.json experiments/results/test/2026-08-26T20-32-48Z_p6-llm_6457ac7/metrics.json
\texttt{tel as tel-national} & 0 & 6 & 0 \\  % results: experiments/results/test/2026-08-26T21-15-35Z_p6-rules_6457ac7/metrics.json experiments/results/test/2026-08-26T21-17-53Z_p6-ngram_6457ac7/metrics.json experiments/results/test/2026-08-26T20-32-48Z_p6-llm_6457ac7/metrics.json
\texttt{tel-national as tel} & 15 & 7 & 14 \\  % results: experiments/results/test/2026-08-26T21-15-35Z_p6-rules_6457ac7/metrics.json experiments/results/test/2026-08-26T21-17-53Z_p6-ngram_6457ac7/metrics.json experiments/results/test/2026-08-26T20-32-48Z_p6-llm_6457ac7/metrics.json
\texttt{username as email} & 0 & 0 & 0 \\  % results: experiments/results/test/2026-08-26T21-15-35Z_p6-rules_6457ac7/metrics.json experiments/results/test/2026-08-26T21-17-53Z_p6-ngram_6457ac7/metrics.json experiments/results/test/2026-08-26T20-32-48Z_p6-llm_6457ac7/metrics.json
\bottomrule
\end{tabular}
\endgroup
\par\smallskip\noindent{\footnotesize Source: \rpath{experiments/results/test/2026-08-26T21-15-35Z_p6-rules_6457ac7/metrics.json}}

\captionof{table}{The confusion pairs named in the specification before any
measurement existed, and how often each actually occurred. Most of them never
did, in either direction, for any engine.}
\label{tab:predicted-pairs}
\end{center}

\section{Form family}

\begin{center}
\small
\begingroup\small\setlength{\tabcolsep}{4pt}
\begin{tabular}{lrrrr}
\toprule
\textbf{Form family} & \textbf{Fields} & \textbf{rules} & \textbf{ngram} & \textbf{llm} \\
\midrule
address & 352 & 0.7222 & 0.4194 & 0.5371 \\  % results: experiments/results/test/2026-08-26T21-15-35Z_p6-rules_6457ac7/metrics.json experiments/results/test/2026-08-26T21-17-53Z_p6-ngram_6457ac7/metrics.json experiments/results/test/2026-08-26T20-32-48Z_p6-llm_6457ac7/metrics.json
checkout & 821 & 0.8414 & 0.6447 & 0.4741 \\  % results: experiments/results/test/2026-08-26T21-15-35Z_p6-rules_6457ac7/metrics.json experiments/results/test/2026-08-26T21-17-53Z_p6-ngram_6457ac7/metrics.json experiments/results/test/2026-08-26T20-32-48Z_p6-llm_6457ac7/metrics.json
login & 216 & 0.6843 & 0.4700 & 0.4476 \\  % results: experiments/results/test/2026-08-26T21-15-35Z_p6-rules_6457ac7/metrics.json experiments/results/test/2026-08-26T21-17-53Z_p6-ngram_6457ac7/metrics.json experiments/results/test/2026-08-26T20-32-48Z_p6-llm_6457ac7/metrics.json
payment & 384 & 0.8570 & 0.5464 & 0.6913 \\  % results: experiments/results/test/2026-08-26T21-15-35Z_p6-rules_6457ac7/metrics.json experiments/results/test/2026-08-26T21-17-53Z_p6-ngram_6457ac7/metrics.json experiments/results/test/2026-08-26T20-32-48Z_p6-llm_6457ac7/metrics.json
signup & 544 & 0.7634 & 0.5521 & 0.5251 \\  % results: experiments/results/test/2026-08-26T21-15-35Z_p6-rules_6457ac7/metrics.json experiments/results/test/2026-08-26T21-17-53Z_p6-ngram_6457ac7/metrics.json experiments/results/test/2026-08-26T20-32-48Z_p6-llm_6457ac7/metrics.json
\bottomrule
\end{tabular}
\endgroup
\par\noindent{\footnotesize Source: \rpath{experiments/results/test/2026-08-26T21-15-35Z_p6-rules_6457ac7/metrics.json}}
\par\noindent{\footnotesize Source: \rpath{experiments/results/test/2026-08-26T21-17-53Z_p6-ngram_6457ac7/metrics.json}}
\par\noindent{\footnotesize Source: \rpath{experiments/results/test/2026-08-26T20-32-48Z_p6-llm_6457ac7/metrics.json}}

\captionof{table}{Macro-F1 by form family.}
\label{tab:by-family}
\end{center}

\section{Label coverage in the corpus}

\begin{center}
\small
\begin{longtable}{lr}
\toprule
\textbf{Label} & \textbf{Keyed fields in the corpus} \\
\midrule
\endfirsthead
\toprule
\textbf{Label} & \textbf{Keyed fields in the corpus} \\
\midrule
\endhead
\bottomrule
\endfoot
\texttt{CC\_EXP\_SPLIT\_MONTH} & 120 \\  % results: corpus/manifest.json
\texttt{CC\_EXP\_SPLIT\_YEAR} & 120 \\  % results: corpus/manifest.json
\texttt{COMPOSITE\_UNSPLIT} & 264 \\  % results: corpus/manifest.json
\texttt{NOT\_AUTOFILLABLE} & 1296 \\  % results: corpus/manifest.json
\texttt{UNKNOWN} & 480 \\  % results: corpus/manifest.json
\texttt{additional-name} & 96 \\  % results: corpus/manifest.json
\texttt{address-level1} & 288 \\  % results: corpus/manifest.json
\texttt{address-level2} & 404 \\  % results: corpus/manifest.json
\texttt{address-line1} & 264 \\  % results: corpus/manifest.json
\texttt{address-line2} & 264 \\  % results: corpus/manifest.json
\texttt{address-line3} & 96 \\  % results: corpus/manifest.json
\texttt{bday} & 96 \\  % results: corpus/manifest.json
\texttt{cc-csc} & 384 \\  % results: corpus/manifest.json
\texttt{cc-exp} & 72 \\  % results: corpus/manifest.json
\texttt{cc-exp-month} & 96 \\  % results: corpus/manifest.json
\texttt{cc-exp-year} & 96 \\  % results: corpus/manifest.json
\texttt{cc-name} & 264 \\  % results: corpus/manifest.json
\texttt{cc-number} & 384 \\  % results: corpus/manifest.json
\texttt{cc-type} & 168 \\  % results: corpus/manifest.json
\texttt{country} & 360 \\  % results: corpus/manifest.json
\texttt{country-name} & 192 \\  % results: corpus/manifest.json
\texttt{current-password} & 168 \\  % results: corpus/manifest.json
\texttt{email} & 600 \\  % results: corpus/manifest.json
\texttt{family-name} & 392 \\  % results: corpus/manifest.json
\texttt{given-name} & 392 \\  % results: corpus/manifest.json
\texttt{honorific-prefix} & 144 \\  % results: corpus/manifest.json
\texttt{honorific-suffix} & 144 \\  % results: corpus/manifest.json
\texttt{name} & 144 \\  % results: corpus/manifest.json
\texttt{new-password} & 360 \\  % results: corpus/manifest.json
\texttt{nickname} & 144 \\  % results: corpus/manifest.json
\texttt{one-time-code} & 240 \\  % results: corpus/manifest.json
\texttt{organization} & 192 \\  % results: corpus/manifest.json
\texttt{postal-code} & 403 \\  % results: corpus/manifest.json
\texttt{sex} & 96 \\  % results: corpus/manifest.json
\texttt{street-address} & 168 \\  % results: corpus/manifest.json
\texttt{tel} & 144 \\  % results: corpus/manifest.json
\texttt{tel-country-code} & 144 \\  % results: corpus/manifest.json
\texttt{tel-extension} & 96 \\  % results: corpus/manifest.json
\texttt{tel-national} & 144 \\  % results: corpus/manifest.json
\texttt{transaction-amount} & 120 \\  % results: corpus/manifest.json
\texttt{url} & 120 \\  % results: corpus/manifest.json
\texttt{username} & 192 \\  % results: corpus/manifest.json
\end{longtable}
\par\smallskip\noindent{\footnotesize Source: \rpath{corpus/manifest.json}}

\captionof{table}{Keyed fields per label across the whole generated corpus. This
is the evidence for the growth rule's clause that every label must be emitted by
at least one answer key, and it is also the honest picture of how uneven the
class distribution is.}
\label{tab:label-coverage}
\end{center}
